\begin{deluxetable}{lc}[tbp]
\centering
\tabletypesize{\footnotesize}
\tablecaption{Properties of the galaxy cluster A370} 
\label{tab:cluster}
\tablehead{ 
 \multicolumn{1}{c}{Parameter} &
 \multicolumn{1}{c}{Value} 
} 
\startdata
ID                                  & A370 \\
Reference center position (J2000.0) & \\
R.A.                        & 02:39:52.9\\
Decl.                       & $-$01:34:36.5\\
X-ray emission centroid (J2000.0)   & \\
R.A.                        & 02:39:53.2\\
Decl.                       & $-$01:34:35.1\\
Redshift                            & $0.375$\\
Velocity dispersion (km~s$^{-1}$)   & $1520\pm 93$\\
X-ray temperature  (keV)            & $8.77^{+0.33}_{-0.34}$\\
Einstein radius ($\arcsec$)         & $33.9\pm 1.1$ for $\zs=2$
\enddata
% core region: 350kpc
\tablecomments{
The optical cluster center is at the midpoint of the two BCGs \citep{Lotz2017hff,Buffalo}. Units of right ascension are hours, minutes, and seconds, and units of declination are degrees, arcminutes, and arcseconds. The cluster velocity dispersion is derived from spectroscopic observations of \citet{Lagattuta2022} in the core region of the cluster \citep[see also][]{Lagattuta2019a370}. The X-ray emission centroid is determined from a 2D $\beta$-model fit to \Chandra X-ray observations (see Section~\ref{sec:xray}). The average temperature of the cluster is measured from the \Chandra X-ray spectrum in the radial range $\in [50,500]\kpch$ centered on the X-ray centroid. The Einstein radius is constrained by detailed strong lens modeling by \citet{Kawamata2018}.
}
\end{deluxetable}

